\documentclass[11pt]{article}

\usepackage[utf8]{inputenc}
\usepackage[T1]{fontenc}
\usepackage[spanish,es-tabla]{babel}
\usepackage[a4paper,margin=2.5cm]{geometry}
\usepackage{hyperref}

\title{Instructivo Completo: Configurar LaTeX + VS Code Desde Cero\\
para Trabajar y Compilar un Proyecto como \texttt{articulo\_rcta.tex}}
\author{Gu\'ia pr\'actica paso a paso}
\date{\today}

\begin{document}
\maketitle
\tableofcontents
\newpage

\section{Objetivo de este instructivo}
Este documento explica, de forma detallada y desde cero, c\'omo dejar operativo un computador reci\'en formateado para editar, compilar y depurar un proyecto LaTeX en VS Code, incluyendo problemas reales que aparecen con frecuencia:

\begin{itemize}
  \item errores de compilaci\'on en LaTeX Workshop,
  \item conflictos de guardado en VS Code (\textit{file is newer}),
  \item salidas de \texttt{latexmk} que dicen que todo est\'a al d\'ia pero muestran errores previos,
  \item y diferencias de codificaci\'on de caracteres.
\end{itemize}

El objetivo final es que puedas modificar \texttt{.tex}, generar tu PDF sin fricci\'on y saber exactamente qu\'e hacer cuando falle algo.

\section{Resultado final esperado}
Al terminar, deber\'ias poder:

\begin{itemize}
  \item abrir una carpeta de proyecto en VS Code,
  \item editar un archivo como \texttt{articulo\_rcta.tex},
  \item compilar desde VS Code o terminal,
  \item ver el PDF actualizado,
  \item y resolver r\'apidamente errores comunes de build.
\end{itemize}

\section{Fase 0: Preparaci\'on inicial del computador}
\subsection{1) Actualizar Windows}
Antes de instalar herramientas de desarrollo:
\begin{itemize}
  \item Ejecuta Windows Update.
  \item Reinicia si te lo pide.
  \item Repite hasta que no queden actualizaciones pendientes importantes.
\end{itemize}

Esto evita problemas de certificados, rutas, componentes faltantes y comportamiento raro en instaladores.

\subsection{2) Crear una carpeta de trabajo}
Por ejemplo:
\begin{verbatim}
D:\pipeline_SVM\article1
\end{verbatim}

Evita rutas demasiado largas o con caracteres especiales no ASCII al principio del proyecto.

\section{Fase 1: Instalar MiKTeX (motor LaTeX)}
\subsection{1) Instalaci\'on}
Instala MiKTeX para tu usuario. Si tienes opci\'on:
\begin{itemize}
  \item arquitectura: 64 bits,
  \item instalaci\'on por usuario (\textit{Just for me}),
  \item paquetes faltantes: instalar autom\'aticamente (\textit{Ask me first} o \textit{Yes}, seg\'un prefieras control o rapidez).
\end{itemize}

\subsection{2) Verificar binarios importantes}
Confirma que existen estos ejecutables (ruta aproximada):
\begin{verbatim}
C:\Users\<TU_USUARIO>\AppData\Local\Programs\MiKTeX\miktex\bin\x64\
  pdflatex.exe
  latexmk.exe
\end{verbatim}

\subsection{3) Actualizar MiKTeX}
Abre MiKTeX Console y ejecuta actualizaciones de paquetes y del sistema MiKTeX. Esto reduce much\'isimo errores de compilaci\'on por paquetes desactualizados.

\section{Fase 2: Instalar VS Code y extensiones}
\subsection{1) Instalar VS Code}
Instala Visual Studio Code (instalaci\'on normal).

\subsection{2) Instalar extensi\'on LaTeX Workshop}
En VS Code:
\begin{itemize}
  \item abre Extensions,
  \item busca \texttt{LaTeX Workshop},
  \item instala la extensi\'on oficial.
\end{itemize}

\section{Fase 3: Configuraci\'on recomendada de VS Code}
\subsection{1) Configuraci\'on base}
En \texttt{settings.json} de usuario, una base razonable es:

\begin{verbatim}
{
  "latex-workshop.latex.autoBuild.run": "onSave",
  "latex-workshop.view.pdf.viewer": "tab",
  "latex-workshop.message.error.show": true,
  "latex-workshop.latex.clean.subfolder.enabled": true,
  "files.autoSave": "off"
}
\end{verbatim}

Notas:
\begin{itemize}
  \item \texttt{onSave} compila al guardar.
  \item Si ves compilaciones en mal momento, puedes cambiar a \texttt{"never"} y compilar manualmente.
  \item \texttt{files.autoSave: off} evita conflictos de guardado innecesarios cuando hay procesos externos tocando archivos.
\end{itemize}

\subsection{2) Validar codificaci\'on del archivo}
Para evitar texto roto (\textit{mojibake}):
\begin{itemize}
  \item En la esquina inferior derecha de VS Code, confirma que el archivo est\'a en \texttt{UTF-8}.
  \item Si no lo est\'a: \texttt{Reopen with Encoding} / \texttt{Save with Encoding} y elige \texttt{UTF-8}.
\end{itemize}

\section{Fase 4: Estructura de proyecto recomendada}
Ejemplo:
\begin{verbatim}
article1/
  articulo_rcta.tex
  Plantilla_CIETA.pdf
  (archivos auxiliares: .aux .log .out ...)
\end{verbatim}

Si usas fondo de plantilla, aseg\'urate de que \texttt{Plantilla\_CIETA.pdf} exista en la misma carpeta cuando el \texttt{.tex} lo referencia.

\section{Fase 5: Flujo de trabajo diario (el m\'as estable)}
\subsection{1) Abrir carpeta}
Abre \texttt{D:\textbackslash pipeline\_SVM\textbackslash article1} como carpeta en VS Code.

\subsection{2) Editar y guardar}
Edita el \texttt{.tex} y guarda.

\subsection{3) Compilar}
Dos maneras:

\textbf{A. Desde VS Code}
\begin{itemize}
  \item Comando: \texttt{LaTeX Workshop: Build LaTeX project}
  \item Atajo usual: \texttt{Ctrl+Alt+B}
\end{itemize}

\textbf{B. Desde terminal PowerShell (ruta completa)}
\begin{verbatim}
& "C:\Users\<TU_USUARIO>\AppData\Local\Programs\MiKTeX\miktex\bin\x64\pdflatex.exe" `
  -interaction=nonstopmode articulo_rcta.tex
\end{verbatim}

Para flujo con \texttt{latexmk}:
\begin{verbatim}
& "C:\Users\<TU_USUARIO>\AppData\Local\Programs\MiKTeX\miktex\bin\x64\latexmk.exe" `
  -pdf -interaction=nonstopmode articulo_rcta.tex
\end{verbatim}

\subsection{4) Ver PDF}
Abre \texttt{articulo\_rcta.pdf} desde el visor de VS Code o desde el explorador.

\section{Fase 6: Soluci\'on de errores frecuentes (muy importante)}
\subsection{Caso A: ``Recipe terminated with error''}
Pasos:
\begin{enumerate}
  \item Abre el \texttt{compiler log}.
  \item Busca la primera l\'inea con \texttt{!} (ese suele ser el error real).
  \item Corrige el \texttt{.tex}.
  \item Limpia auxiliares y recompila.
\end{enumerate}

\subsection{Caso B: ``Failed to save ... file is newer''}
Esto pasa cuando el archivo en disco cambi\'o y tu pesta\~na tiene una versi\'on anterior.
\begin{enumerate}
  \item Pulsa \texttt{Compare}.
  \item Fusiona o conserva la versi\'on correcta.
  \item Guarda de nuevo.
  \item Recompila.
\end{enumerate}

Consejo: evita editar el mismo archivo simult\'aneamente con varios programas.

\subsection{Caso C: \texttt{latexmk} dice ``up-to-date'' pero muestra error previo}
Ese estado es com\'un. Soluci\'on robusta:
\begin{verbatim}
& "C:\Users\<TU_USUARIO>\AppData\Local\Programs\MiKTeX\miktex\bin\x64\latexmk.exe" `
  -C articulo_rcta.tex

& "C:\Users\<TU_USUARIO>\AppData\Local\Programs\MiKTeX\miktex\bin\x64\latexmk.exe" `
  -pdf -interaction=nonstopmode -g articulo_rcta.tex
\end{verbatim}

\subsection{Caso D: Mensajes de code page en consola}
L\'ineas como:
\begin{verbatim}
Initial Win CP ...
I changed them all to CP1252
\end{verbatim}
normalmente son informativas, no son el error en s\'i.

\subsection{Caso E: Error por porcentajes con espa\~nol}
En algunos entornos con \texttt{babel-spanish}, el patr\'on:
\begin{verbatim}
\,\%
\end{verbatim}
puede disparar errores del tipo \texttt{Incompatible glue units}. Soluci\'on pr\'actica:
\begin{itemize}
  \item usar simplemente \texttt{\textbackslash\%} en el texto normal,
  \item y evitar combinaciones extra\~nas de espacio fino antes del porcentaje.
\end{itemize}

\section{Fase 7: Comandos de recuperaci\'on r\'apida}
Cuando ``todo est\'a raro'', ejecuta este bloque:

\begin{verbatim}
& "C:\Users\<TU_USUARIO>\AppData\Local\Programs\MiKTeX\miktex\bin\x64\latexmk.exe" `
  -C articulo_rcta.tex

& "C:\Users\<TU_USUARIO>\AppData\Local\Programs\MiKTeX\miktex\bin\x64\latexmk.exe" `
  -pdf -interaction=nonstopmode -g articulo_rcta.tex
\end{verbatim}

Si no usas \texttt{latexmk}, compila dos veces con \texttt{pdflatex}:
\begin{verbatim}
& "C:\Users\<TU_USUARIO>\AppData\Local\Programs\MiKTeX\miktex\bin\x64\pdflatex.exe" `
  -interaction=nonstopmode articulo_rcta.tex

& "C:\Users\<TU_USUARIO>\AppData\Local\Programs\MiKTeX\miktex\bin\x64\pdflatex.exe" `
  -interaction=nonstopmode articulo_rcta.tex
\end{verbatim}

\section{Fase 8: Buenas pr\'acticas para no volver al caos}
\begin{itemize}
  \item Usa control de versiones (Git) desde el d\'ia 1.
  \item Haz commits peque\~nos despu\'es de cambios que compilan.
  \item Mant\'en encoding UTF-8 siempre.
  \item Evita ``copiar y pegar'' texto desde fuentes que cambian codificaci\'on.
  \item Limpia auxiliares cuando cambies estructura grande de secciones/referencias.
  \item Si VS Code muestra conflicto de guardado, primero compara y luego guarda.
\end{itemize}

\section{Checklist final (10 minutos)}
\begin{enumerate}
  \item MiKTeX instalado y actualizado.
  \item VS Code + LaTeX Workshop instalados.
  \item Archivo en UTF-8.
  \item Build ejecuta sin errores fatales.
  \item PDF se genera y abre.
  \item Ya sabes correr limpieza + rebuild forzado.
\end{enumerate}

\section{Cierre}
Si sigues este instructivo, vas a poder reproducir un entorno estable desde un equipo reci\'en formateado y trabajar de forma predecible. La clave no es ``que nunca falle'', sino tener un procedimiento claro para recuperar el estado de compilaci\'on en menos de cinco minutos.

\vspace{1em}
\noindent
Sugerencia final: guarda este instructivo dentro del proyecto y actual\'izalo cada vez que descubras un nuevo error recurrente. Convertir problemas reales en pasos repetibles es lo que vuelve robusto tu flujo de trabajo.

\end{document}

